\documentclass[reqno]{amsart}
\usepackage[T1]{fontenc}
\usepackage{amsmath,amssymb,amsthm}
\usepackage{xr-hyper}
\usepackage[hidelinks]{hyperref}
\externaldocument[ps-]{20261125T010500--poincare-series-and-near-orthogonality__workbook}
\externaldocument[tn-]{20261125T002100--definite-theta-lifts-and-norms__workbook}
\externaldocument[bp-]{20261126T020210--bessel-poincare-series-and-their-gram-operators__workbook}
\externaldocument[ss-]{20261126T030400--the-seesaw-identity__workbook}
\externaldocument[ob-]{20261126T040100--spaces-forms-and-twists__workbook}

% Chapter number; \ref{chap} gives it.
\newcounter{chap}
% Chapter numbers, generated from 20261125T012000--grh-workbook__guide.org
% by scripts/make-chapters.py.  Do not edit: change the order of the
% index and rerun the script (./preview-site.sh and CI do this).
%
% Usage in a sheet, after \newcounter{chap}:
%   % Chapter numbers, generated from 20261125T012000--grh-workbook__guide.org
% by scripts/make-chapters.py.  Do not edit: change the order of the
% index and rerun the script (./preview-site.sh and CI do this).
%
% Usage in a sheet, after \newcounter{chap}:
%   % Chapter numbers, generated from 20261125T012000--grh-workbook__guide.org
% by scripts/make-chapters.py.  Do not edit: change the order of the
% index and rerun the script (./preview-site.sh and CI do this).
%
% Usage in a sheet, after \newcounter{chap}:
%   \input{chapters}\chapnumber{<sheet ID>}
\makeatletter
\providecommand\chapters@set[2]{\@namedef{chapters@#1}{#2}}
\chapters@set{20261125T001000}{1}% zeros-to-ramp-sums
\chapters@set{20261125T001100}{2}% inert-primes-and-estimate-b
\chapters@set{20261125T005000}{3}% galois-fourier-and-averaging
\chapters@set{20261125T005500}{4}% columns-from-a-sphere
\chapters@set{20261125T000000}{5}% slope-method
\chapters@set{20261126T030100}{6}% a-twist-for-each-prime
\chapters@set{20261125T003000}{7}% the-chain-and-the-square-root
\chapters@set{20261125T010000}{8}% rank-lemma-and-zero-estimate
\chapters@set{20261125T010500}{9}% poincare-series-and-near-orthogonality
\chapters@set{20261124T223000}{10}% elliptic-curves-for-the-proof
\chapters@set{20261125T009500}{11}% fock-model-and-theta-intertwiners
\chapters@set{20261124T223100}{12}% complex-multiplication-and-inert-reduction
\chapters@set{20261125T008000}{13}% one-cm-determinant
\chapters@set{20261124T214500}{14}% supersingular-theta-divisibility
\chapters@set{20261124T214600}{15}% the-isotropic-plane
\chapters@set{20261125T013700}{16}% pairings-on-a-formal-group-of-height-two
\chapters@set{20261125T013500}{17}% hecke-nondegeneracy-on-the-tree
\chapters@set{20261126T030200}{18}% the-bound-is-attained
\chapters@set{20261125T009000}{19}% kronecker-limit-and-chowla-selberg
\chapters@set{20261125T002000}{20}% definite-theta-lifts
\chapters@set{20261125T015500}{21}% siegel-weil-in-the-smallest-case
\chapters@set{20261125T002100}{22}% definite-theta-lifts-and-norms
\chapters@set{20261125T016500}{23}% hecke-operators-and-theta-lifts
\chapters@set{20261125T016700}{24}% divisibility-along-the-chain
\chapters@set{20261125T004000}{25}% denominators-and-sha
\chapters@set{20261125T004100}{26}% tate-twists-and-sha
\chapters@set{20261126T020430}{27}% the-minus-part-by-gauss-sums
\chapters@set{20261125T014100}{28}% congruence-modules
\chapters@set{20261125T014500}{29}% eisenstein-congruences-and-selmer-classes
\chapters@set{20261125T014700}{30}% pseudocharacters-and-the-length-bound
\chapters@set{20261125T011000}{31}% galois-representations-black-box
\chapters@set{20261126T030300}{32}% the-data
\chapters@set{20261126T020210}{33}% bessel-poincare-series-and-their-gram-operators
\chapters@set{20261126T030400}{34}% the-seesaw-identity
\chapters@set{20261126T030500}{35}% divisibility-of-the-composite-lift
\chapters@set{20261126T030600}{36}% theorem-a-short-forms
\chapters@set{20261126T030700}{37}% the-determinant-of-a-chain
\chapters@set{20261126T030800}{38}% the-cost-of-one-lift
\chapters@set{20261126T030900}{39}% denominators-of-projections-to-lifts
\chapters@set{20261126T031000}{40}% twists-with-few-congruences
\chapters@set{20261126T031100}{41}% theorem-c-and-the-main-inequality
\providecommand\chapnumber[1]{%
  \@ifundefined{chapters@#1}%
    {\PackageError{chapters}{Sheet #1 is not listed in the index}%
      {List it in 20261125T012000--grh-workbook__guide.org\MessageBreak
       and run scripts/make-chapters.py.}}%
    {\edef\thechap{\@nameuse{chapters@#1}}}}
\makeatother
\chapnumber{<sheet ID>}
\makeatletter
\providecommand\chapters@set[2]{\@namedef{chapters@#1}{#2}}
\chapters@set{20261125T001000}{1}% zeros-to-ramp-sums
\chapters@set{20261125T001100}{2}% inert-primes-and-estimate-b
\chapters@set{20261125T005000}{3}% galois-fourier-and-averaging
\chapters@set{20261125T005500}{4}% columns-from-a-sphere
\chapters@set{20261125T000000}{5}% slope-method
\chapters@set{20261126T030100}{6}% a-twist-for-each-prime
\chapters@set{20261125T003000}{7}% the-chain-and-the-square-root
\chapters@set{20261125T010000}{8}% rank-lemma-and-zero-estimate
\chapters@set{20261125T010500}{9}% poincare-series-and-near-orthogonality
\chapters@set{20261124T223000}{10}% elliptic-curves-for-the-proof
\chapters@set{20261125T009500}{11}% fock-model-and-theta-intertwiners
\chapters@set{20261124T223100}{12}% complex-multiplication-and-inert-reduction
\chapters@set{20261125T008000}{13}% one-cm-determinant
\chapters@set{20261124T214500}{14}% supersingular-theta-divisibility
\chapters@set{20261124T214600}{15}% the-isotropic-plane
\chapters@set{20261125T013700}{16}% pairings-on-a-formal-group-of-height-two
\chapters@set{20261125T013500}{17}% hecke-nondegeneracy-on-the-tree
\chapters@set{20261126T030200}{18}% the-bound-is-attained
\chapters@set{20261125T009000}{19}% kronecker-limit-and-chowla-selberg
\chapters@set{20261125T002000}{20}% definite-theta-lifts
\chapters@set{20261125T015500}{21}% siegel-weil-in-the-smallest-case
\chapters@set{20261125T002100}{22}% definite-theta-lifts-and-norms
\chapters@set{20261125T016500}{23}% hecke-operators-and-theta-lifts
\chapters@set{20261125T016700}{24}% divisibility-along-the-chain
\chapters@set{20261125T004000}{25}% denominators-and-sha
\chapters@set{20261125T004100}{26}% tate-twists-and-sha
\chapters@set{20261126T020430}{27}% the-minus-part-by-gauss-sums
\chapters@set{20261125T014100}{28}% congruence-modules
\chapters@set{20261125T014500}{29}% eisenstein-congruences-and-selmer-classes
\chapters@set{20261125T014700}{30}% pseudocharacters-and-the-length-bound
\chapters@set{20261125T011000}{31}% galois-representations-black-box
\chapters@set{20261126T030300}{32}% the-data
\chapters@set{20261126T020210}{33}% bessel-poincare-series-and-their-gram-operators
\chapters@set{20261126T030400}{34}% the-seesaw-identity
\chapters@set{20261126T030500}{35}% divisibility-of-the-composite-lift
\chapters@set{20261126T030600}{36}% theorem-a-short-forms
\chapters@set{20261126T030700}{37}% the-determinant-of-a-chain
\chapters@set{20261126T030800}{38}% the-cost-of-one-lift
\chapters@set{20261126T030900}{39}% denominators-of-projections-to-lifts
\chapters@set{20261126T031000}{40}% twists-with-few-congruences
\chapters@set{20261126T031100}{41}% theorem-c-and-the-main-inequality
\providecommand\chapnumber[1]{%
  \@ifundefined{chapters@#1}%
    {\PackageError{chapters}{Sheet #1 is not listed in the index}%
      {List it in 20261125T012000--grh-workbook__guide.org\MessageBreak
       and run scripts/make-chapters.py.}}%
    {\edef\thechap{\@nameuse{chapters@#1}}}}
\makeatother
\chapnumber{<sheet ID>}
\makeatletter
\providecommand\chapters@set[2]{\@namedef{chapters@#1}{#2}}
\chapters@set{20261125T001000}{1}% zeros-to-ramp-sums
\chapters@set{20261125T001100}{2}% inert-primes-and-estimate-b
\chapters@set{20261125T005000}{3}% galois-fourier-and-averaging
\chapters@set{20261125T005500}{4}% columns-from-a-sphere
\chapters@set{20261125T000000}{5}% slope-method
\chapters@set{20261126T030100}{6}% a-twist-for-each-prime
\chapters@set{20261125T003000}{7}% the-chain-and-the-square-root
\chapters@set{20261125T010000}{8}% rank-lemma-and-zero-estimate
\chapters@set{20261125T010500}{9}% poincare-series-and-near-orthogonality
\chapters@set{20261124T223000}{10}% elliptic-curves-for-the-proof
\chapters@set{20261125T009500}{11}% fock-model-and-theta-intertwiners
\chapters@set{20261124T223100}{12}% complex-multiplication-and-inert-reduction
\chapters@set{20261125T008000}{13}% one-cm-determinant
\chapters@set{20261124T214500}{14}% supersingular-theta-divisibility
\chapters@set{20261124T214600}{15}% the-isotropic-plane
\chapters@set{20261125T013700}{16}% pairings-on-a-formal-group-of-height-two
\chapters@set{20261125T013500}{17}% hecke-nondegeneracy-on-the-tree
\chapters@set{20261126T030200}{18}% the-bound-is-attained
\chapters@set{20261125T009000}{19}% kronecker-limit-and-chowla-selberg
\chapters@set{20261125T002000}{20}% definite-theta-lifts
\chapters@set{20261125T015500}{21}% siegel-weil-in-the-smallest-case
\chapters@set{20261125T002100}{22}% definite-theta-lifts-and-norms
\chapters@set{20261125T016500}{23}% hecke-operators-and-theta-lifts
\chapters@set{20261125T016700}{24}% divisibility-along-the-chain
\chapters@set{20261125T004000}{25}% denominators-and-sha
\chapters@set{20261125T004100}{26}% tate-twists-and-sha
\chapters@set{20261126T020430}{27}% the-minus-part-by-gauss-sums
\chapters@set{20261125T014100}{28}% congruence-modules
\chapters@set{20261125T014500}{29}% eisenstein-congruences-and-selmer-classes
\chapters@set{20261125T014700}{30}% pseudocharacters-and-the-length-bound
\chapters@set{20261125T011000}{31}% galois-representations-black-box
\chapters@set{20261126T030300}{32}% the-data
\chapters@set{20261126T020210}{33}% bessel-poincare-series-and-their-gram-operators
\chapters@set{20261126T030400}{34}% the-seesaw-identity
\chapters@set{20261126T030500}{35}% divisibility-of-the-composite-lift
\chapters@set{20261126T030600}{36}% theorem-a-short-forms
\chapters@set{20261126T030700}{37}% the-determinant-of-a-chain
\chapters@set{20261126T030800}{38}% the-cost-of-one-lift
\chapters@set{20261126T030900}{39}% denominators-of-projections-to-lifts
\chapters@set{20261126T031000}{40}% twists-with-few-congruences
\chapters@set{20261126T031100}{41}% theorem-c-and-the-main-inequality
\providecommand\chapnumber[1]{%
  \@ifundefined{chapters@#1}%
    {\PackageError{chapters}{Sheet #1 is not listed in the index}%
      {List it in 20261125T012000--grh-workbook__guide.org\MessageBreak
       and run scripts/make-chapters.py.}}%
    {\edef\thechap{\@nameuse{chapters@#1}}}}
\makeatother
\chapnumber{20261126T051200}
\renewcommand{\thesection}{\thechap.\arabic{section}}

\theoremstyle{definition}
\newtheorem{problem}{Problem}[section]
\newtheorem{definition}{Definition}[section]
\theoremstyle{plain}
\newtheorem{theorem}{Theorem}[section]
\theoremstyle{remark}
\newtheorem*{leanremark}{In the Lean formalization}

\newcommand{\Z}{\mathbb{Z}}
\newcommand{\C}{\mathbb{C}}
\newcommand{\F}{\mathbb{F}}
\newcommand{\A}{\mathbb{A}}
\newcommand{\ii}{\mathrm{i}}
\newcommand{\Uc}{\mathcal{U}}
\newcommand{\Wc}{\mathcal{W}}
\newcommand{\Cc}{\mathcal{C}}
\newcommand{\Sc}{\mathcal{S}}
\newcommand{\WW}{\mathbb{W}}
\newcommand{\st}{\mathrm{st}}
\newcommand{\CT}{\mathrm{CT}}
\newcommand{\Heis}{\mathrm{Heis}}
\newcommand{\Stab}{\operatorname{Stab}}
\newcommand{\Hom}{\operatorname{Hom}}
\newcommand{\tr}{\operatorname{tr}}
\newcommand{\Tr}{\operatorname{Tr}}
\newcommand{\len}{\operatorname{len}}
\newcommand{\rep}{\mathrm{rep}}
\newcommand{\tl}[1]{\texttt{#1}}
\newcommand{\glsrc}[3]{\href{https://github.com/lsernau/grh-lean/blob/72fead1/grh/src/#1\#L#2}{#3}}

\setlength{\emergencystretch}{3em}

\title{Transferring a Bessel period through a theta lift}
\author{}
\date{November 26, 2026}

\begin{document}
\maketitle
\refstepcounter{chap}\label{chap}

\paragraph{Prerequisites.}
Hermitian forms over fields; Witt's extension theorem \cite{Sch}.  From Chapter~\ref{bp-chap}: the Bessel subgroup \(R=\Uc\rtimes U(X)\) and an example of its character (Section~\ref{bp-sec:subgroup}, Problems~\ref{bp-pr:lambda} and~\ref{bp-pr:normalize}), and the finite unfolding of Problem~\ref{bp-pr:unfold}.  From Chapter~\ref{ps-chap}: the unfolding of Problem~\ref{ps-pr:unfold}.  From Chapter~\ref{tn-chap}: Problem~\ref{tn-pr:adjoint} and Section~\ref{tn-sec:seesaw}.  From Chapter~\ref{ss-chap}: Theorem~\ref{ss-thm:seesaw} and Problem~\ref{ss-pr:transpose}.  The Weil representation of a unitary dual pair and the Heisenberg representation under it, at the level of their definitions \cite{We}, \cite[\S1]{Ku}.  Matrices act on columns; the matrix of a map \(V\to V'\) has its rows indexed by the basis of \(V'\).  In this sheet \(b\), \(d\), \(s\) are local letters (not the scale, the prime and the tower exponent).

\medskip
The proof of Theorem~\ref{ss-thm:seesaw} pairs both sides with a form on the last space of the lower chain and moves the lifts onto that form by transposition, as in Problem~\ref{ss-pr:transpose}; what is left is a Bessel period of a theta function, integrated against the form.  How does such a period pass through one theta lift?  For the weights \(0\), the theta function expands in constant terms indexed by the classes of rational maps \(V\to V'\) modulo maps of degree \(\ge0\) (degrees are defined in Section~\ref{sec:strings}, and the expansion is stated in Section~\ref{sec:transfer}).  In the period against a Bessel function only the \(U(V')(F)\)-translates of the set of classes agreeing with the string map contribute; the rational Bessel subgroup of the second space preserves this set, and each translate contributes one value of the transferred function.  We first see the mechanism over \(\F_3\) and \(\F_9\).


\section{Theta against a character on a plane over \texorpdfstring{\(\F_3\)}{F3}}\label{sec:plane}

Let \(Q(v)=v_1^2+v_2^2\) on \(\F_3^2\), let \(\zeta_3:=e^{2\pi\ii/3}\), and for a function \(\phi\colon\F_3^2\to\C\) and \(b\in\F_3\) put \((n(b)\phi)(v):=\zeta_3^{bQ(v)}\phi(v)\) and \(\theta(\phi):=\sum_v\phi(v)\).  The operators \(n(b)\) commute with \((h\phi)(v):=\phi(h^{-1}v)\) for \(h\) in the orthogonal group \(O(Q)\).  In this section \(\phi(v):=1+[v_1]+2[v_2]\), with \([\,\cdot\,]\in\{0,1,2\}\) the least residue.

\begin{problem}\label{pr:buckets}
The sums \(\sum_{Q(v)=1+j}\phi(v)\) are \(13\), \(22\), \(1\) for \(j=0,1,2\).
\end{problem}

\begin{problem}\label{pr:coefficient}
For every \(\phi\), \(\frac13\sum_{b\in\F_3}\zeta_3^{-b}\,\theta(n(b)\phi)=\sum_{Q(v)=1}\phi(v)\).  For the \(\phi\) above it is \(13\).
\end{problem}

The left side is a period of \(b\mapsto\theta(n(b)\phi)\) against a character; only the \(v\) with \(Q(v)=1\) survive it.

\begin{problem}\label{pr:plane-orbit}
\(O(Q)\) has \(8\) elements, the stabilizer of \(v_0:=(1,0)\) has \(2\), and the orbit of \(v_0\) is \(\{(\pm1,0),(0,\pm1)\}=\{Q=1\}\).
\end{problem}

\begin{problem}\label{pr:plane-unfold}
Let \(\mathrm{cf}(\phi)\) be the coefficient of Problem~\ref{pr:coefficient} and \(\Stab\) the stabilizer of \(v_0\).  For \(h\in O(Q)\), \(\mathrm{cf}(h\phi)=\sum_{q\in O(Q)/\Stab}\phi(h^{-1}qv_0)\); for \(h=1\) this is \(\frac12\sum_{\gamma\in O(Q)}\phi(\gamma v_0)=13\).  For every function \(z\) on \(O(Q)\),
\[
  \sum_{h\in O(Q)}z(h)\,\mathrm{cf}(h\phi)=\sum_{h\in O(Q)}\Bigl(\frac{1}{\lvert\Stab\rvert}\sum_{\gamma\in O(Q)}z(\gamma h)\Bigr)\phi(h^{-1}v_0).
\]
\end{problem}

At fixed \(h\), the coefficient \(\mathrm{cf}(h\phi)\) is the sum over \(O(Q)/\Stab\) of the function \(k\mapsto\phi(k^{-1}v_0)\) evaluated at \(k=q^{-1}h\); this is the finite model of the transfer.  In the adelic period the unipotent group can be nonabelian, and the proof uses two choices of weights.


\section{Strings, grades and the string map}\label{sec:strings}

For the matrix definitions and the finite examples, let \(E\) be a field with an involution \(x\mapsto\bar x\) and fixed field \(F\), and let \(N\in F^\times\); all bases are finite.

\begin{definition}\label{def:string}
A \emph{space with a string} for \(N\) is a hermitian space \(V\) over \(E\) with a basis consisting of an orthonormal basis of a subspace \(X\) and vectors \(f_0,\dots,f_{s-1}\) orthogonal to \(X\), and with a constant \(c_V\in F^\times\), such that \(f_i\) pairs only with \(f_{s-1-i}\), to \(1\), or to \(-N\) if \(i=s-1-i\).  We write \(J_V\) for the Gram matrix.  \(X\) is the \emph{definite part}, \(f_0,\dots,f_{s-1}\) is the \emph{string}, and \(\len V:=s\) is its \emph{length}.
\end{definition}

\begin{definition}\label{def:degree}
The \emph{grade} of a vector of the basis of \(X\) is \(0\), and that of \(f_i\) is \(s-1-2i\).  Let \(V'\) be a second space with a string, and let \(d\) be integer weights on the basis of \(V'\).  For basis vectors \(e\) of \(V\) and \(e'\) of \(V'\), the entry of a map \(w\colon V\to V'\) in the row of \(e'\) and the column of \(e\) has \emph{degree} \(d(e')-\operatorname{grade}(e)\), and \(w\) has degree \(\ge i\) if all its entries of degree \(<i\) vanish.  For maps \(V\to V\) the weights are the grades.
\end{definition}

For the shape of Section~\ref{bp-sec:subgroup} the string is \(v_1,\dots,v_\varrho,z,v'_\varrho,\dots,v'_1\); \(\Uc\) is the group of unitary \(n\) with \(n-1\) raising grades, and \(R\) that of the unitary \(g\) with \(gf_i-f_i\) in the span of the string vectors of larger grade for every \(f_i\) of grade \(\ge0\).  In the adelic setting of the proof these conditions define \(\Uc(\A_F)\) and \(R(\A_F)\) in the string basis.  We use the same matrix conditions over finite fields.

\begin{definition}\label{def:stringmap}
Let \(\len V'=\len V-1\), or \(\len V=\len V'=0\).  The \emph{string map} \(\st\colon V\to V'\) kills the definite part and sends \(f_i\) to \(f'_i\) times a scalar, and to \(0\) if \(V'\) has no vector \(f'_i\).  If \(\len V\) is odd, the scalar is \(1,\ -Nc_V,\ c_V\) for \(f_i\) of positive, zero, negative grade; if \(\len V\) is even, it is \(c_V,\ 1\) for positive, negative grade.  For a map \(w\colon V\to V'\) put \(w^*:=J_V^{-1}\bar w^{t}J_{V'}\).
\end{definition}

For the grades of \(V\) and \(V'\) the string map has degree \(\ge-1\); in the chain of the data it is \(\mathrm{st}_j\colon Y_{j-1}\to Y_j\).

\emph{The hook.}  In the finite examples of this section and the next, \(E=\F_9=\F_3[\ii]\) with \(\ii^2=-1\) and \(\bar\ii=-\ii\), \(N=1\), and \(V\), \(V'\) have no definite part, strings of lengths \(3\) and \(2\), and constants \(c_V=c_{V'}=1\): \(J_V=\operatorname{antidiag}(1,-1,1)\), \(J_{V'}=\operatorname{antidiag}(1,1)\), the grades are \(2,0,-2\) on \(V\) and \(1,-1\) on \(V'\), and \(\st=\left(\begin{smallmatrix}1&0&0\\0&-1&0\end{smallmatrix}\right)\).

\begin{problem}\label{pr:hook-unipotent}
The unitary upper unitriangular matrices for \(J_V\) are the \(27\) matrices \(u=\left(\begin{smallmatrix}1&a&b\\0&1&\bar a\\0&0&1\end{smallmatrix}\right)\) with \(b+\bar b=a\bar a\).  They form \(\Uc\), and \(u-1\) has degree \(\ge2\) for each of them.
\end{problem}

\begin{problem}\label{pr:hook-phase}
\(\tr(\st^*\st\,u)=a-\bar a\), and for \(w=(v,0,0)\) with first column \(v\in\F_9^2\), \(\tr(w^*wu)=(\bar v^{t}J_{V'}v)\,b\).
\end{problem}


\section{The survival test}\label{sec:survival}

Over the adeles, let \(E\) be a CM field with the CM type and Weil data of \cite{GL}, \(F\) its maximal totally real subfield and \(\A_E\) its adeles; \(\psi_E\) is the standard additive character of \(\A_E/E\), and \(j_E\in E^\times\), with \(\bar j_E=-j_E\), is the element given by the Weil data.  In the proof the character used on the unitary \(n\) with \(n-1\) of degree \(\ge2\) is \(\psi_\Uc(n)=\psi_E\bigl(\tfrac12j_E\tr(\st^*\st\,n)\bigr)\); the matrix \(\st^*\st\) is the shift of the string of \(V\) and depends on \(c_V\), and \(\psi_\Uc(nn')=\psi_\Uc(n)\psi_\Uc(n')\) for two such \(n,n'\).  Over \(\F_9\) we use the same formulas with \(\psi_E(x)=\zeta_3^{\Tr_{\F_9/\F_3}x}\) and \(j_E=\ii\) (and \(\frac12=-1\) in \(\F_3\)), and \(u\) runs over \(\Uc(\F_3)\).

\begin{definition}\label{def:survive}
Let \(d\) be weights on the basis of \(V'\).  A rational map \(w\colon V\to V'\) \emph{survives} for \(d\) if
\[
  \psi_E\bigl(\tfrac12j_E\tr(\st^*\st\,u)\bigr)=\psi_E\bigl(\tfrac12j_E\tr(w^*w\,u)\bigr)
\]
for every \(u\in\Uc(\A_F)\) such that \(u-1\) has degree \(\ge2\) and \(w(u-1)\) has degree \(\ge1\) for the weights \(d\).  A map \(w\) \emph{agrees with} \(\st\) if \(w\) and \(\st\) have the same entry in the row of \(e'\) and the column of \(e\) whenever \(\operatorname{grade}(e)>0\) and \(\operatorname{grade}(e')<\operatorname{grade}(e)\).
\end{definition}

The test is applied to classes of rational maps modulo the maps of degree \(\ge0\), through a chosen representative of each class; a class is determined by the entries of negative degree.  For the weights \(0\) these are the columns of positive grade.

\begin{problem}\label{pr:hook-survivors}
In the hook, take the weights \((0,0)\).  The classes are represented by the \(81\) maps \(w=(v,0,0)\), and exactly \(32\) of them survive: the \(v\ne0\) with \(\bar v^{t}J_{V'}v=0\).
\end{problem}

\begin{problem}\label{pr:hook-orbit}
The unitary group \(U(V')(\F_3)=\{g\in\mathrm{GL}_2(\F_9):\bar g^{t}J_{V'}g=J_{V'}\}\) has \(96\) elements, and their first columns are the \(32\) surviving vectors of Problem~\ref{pr:hook-survivors}.  The stabilizer of \((1,0)^{t}\) has order \(3\) and is the Bessel subgroup \(R'\) of \(V'\).
\end{problem}

\begin{problem}\label{pr:hook-second}
Take the weights \((1,-1)\), the grades of \(V'\).  The classes that agree with \(\st\) are represented by the nine maps \(\left(\begin{smallmatrix}1&0&0\\0&c&0\end{smallmatrix}\right)\), \(c\in\F_9\), and exactly \(c=-1\) survives: the class of \(\st\).
\end{problem}

The hook illustrates the two choices of weights in the proof.  For the weights \(0\) the survivors are the translates by \(U(V')(\F_3)\) of the one class that agrees with \(\st\), \(v=(1,0)^{t}\), in bijection with \(U(V')(\F_3)/R'\) (\(96/3=32\)); for the grades of \(V'\) only the class of \(\st\) survives among the agreeing classes.  Here is a finite model of how a failing map is killed.

\begin{problem}\label{pr:twisted-mean}
Let \(H\) be a finite group and \(\mathfrak f\colon H\to\C\) a function with \(\mathfrak f(ux)=\zeta\mathfrak f(x)\) for all \(x\), for one \(u\in H\) and one \(\zeta\ne1\).  Then \(\sum_{x\in H}\mathfrak f(x)=0\).  If \(\mathfrak f\) is also left invariant under a subgroup \(S\), then \(\sum_{x\in S\backslash H}\mathfrak f(x)=0\).
\end{problem}


\section{Constant terms and the transfer}\label{sec:transfer}

Let \(V\), \(V'\) be spaces with strings for the same \(N\) over \(E\), with \(\len V'=\len V-1\) or \(\len V=\len V'=0\), each with an adelic frame (Section~\ref{ob-sec:spaces}) in which its Gram matrix is the same diagonal matrix of signs at every infinite place and the identity at the finite places; matrices are written in the bases of Definition~\ref{def:string}.  Let \(R\), \(R'\) be the Bessel subgroups.  Write \(\A=\A_F\) for the groups and \(\A_E\) for the matrices.  Put \([R]:=R(F)\backslash R(\A)\) and \([U(V')]:=U(V')(F)\backslash U(V')(\A)\), with the quotient measures of fixed Haar measures; \([R]\) is compact.  Let \(\WW:=\Hom(V,V')\), with the pairing \((y,y')\mapsto\tr(y'^*y)\).  For each finite index set fix a continuous unitary character of the norm-one adeles of \(E\) trivial on the norm-one elements of \(E\); for the index sets of the bases of \(V\) and \(V'\) these give the Weil representation \(\omega\) of \(U(V)(\A)\times U(V')(\A)\) on the Banach space \(\Sc\) of data of \cite{GL} (see \cite{We}, \cite[\S1]{Ku}), with the Heisenberg operators \(\Heis(y)\), \(y\in\WW(\A_E)\), and the theta functional \(\theta\colon\Sc\to\C\).  We use: \(\omega(g,h)\Heis(y)\omega(g,h)^{-1}=\Heis(hyg^{-1})\); \(\theta\circ\omega(\gamma,\gamma')=\theta\) for rational \(\gamma,\gamma'\); \(y\mapsto\theta(\Heis(y)\phi)\) is invariant under rational translations on every rational subspace of \(\WW\) on which the pairing vanishes.  \emph{Status}: classical.

Let \(d\) be weights on the basis of \(V'\) that change sign under the involution of the basis fixing the definite part and exchanging \(f'_i\) and \(f'_{s'-1-i}\) (\(s'=\len V'\)), for instance \(d=0\) or the grading, and let \(\WW^d_{>0}\) be the maps whose entries of degree \(\le0\) vanish.  The pairing vanishes on \(\WW^d_{>0}\).  A map \(V'\to V'\) has degree \(\ge0\) for \(d\) if its entry in the row of \(e'\) and column of \(e''\) vanishes whenever \(d(e')<d(e'')\).

\begin{definition}\label{def:ct}
For a datum \(\phi\) and a map \(w\in\WW(\A_E)\), define the \emph{constant term} of \(\phi\) at \(w\) by the mean
\[
  \CT^d_w(\phi):=\int_{\WW^d_{>0}(E)\backslash\WW^d_{>0}(\A_E)}\theta\bigl(\Heis(y)\Heis(w)\phi\bigr)\,dy
\]
for the probability Haar measure.  We write \(\CT_w:=\CT^d_w\) when \(d\) is the grading of \(V'\).
\end{definition}

In the proof, for rational \(w\), \(\CT^d_w\) depends only on \(w\) modulo rational maps of degree \(\ge0\).  For the fixed spaces and weights \(d\) there is a constant \(C_d\ge0\) such that for every datum \(\phi\),
\[
  \theta(\phi)=\sum_w\CT^d_w(\phi),\qquad \sum_w\lvert\CT^d_w(\phi)\rvert\le C_d\lVert \phi\rVert,
\]
the sums over the classes of rational maps modulo maps of degree \(\ge0\), with the norm of \(\Sc\).  This is the Fourier expansion along \(\WW^d_{>0}\) in a mixed model; compare \cite{We}.  \emph{Status}: classical; the norm estimate is used in the form proved in the argument.

\begin{definition}\label{def:transfer}
Let \(\phi\in\Sc\), and let \(\Wc\) be a continuous function on \(R(\A)\), left invariant under \(R(F)\), with \(\Wc(nr)=\psi_\Uc(n)\Wc(r)\) for all \(n\in R(\A)\) such that \(n-1\) has degree \(\ge2\).  Define the \emph{transfer} of \(\Wc\) on \(U(V')(\A)\) by
\[
  \Wc'(h):=\int_{[R]}\Wc(r)\,\CT_\st\bigl(\omega(r,h)\phi\bigr)\,dr .
\]
\end{definition}

In the proof \(\Wc'\) is continuous and left \(R'(F)\)-invariant.  For each class \(w\) choose a rational representative \(\tilde w\), and put \(\mathcal Q^d_w(h,r):=\Wc(r)\,\CT^d_{\tilde w}(\omega(r,h)\phi)\); a class survives, or agrees with \(\st\), if \(\tilde w\) does.  In the proof, each \(\mathcal Q^d_w\) is jointly continuous in \((h,r)\), and there is a constant \(A_d\ge0\) such that \(\sum_{w\in\Sigma}\lvert\mathcal Q^d_w(h,r)\rvert\le A_d\lVert\omega(1,\delta h)\rVert\) for every \(h\in U(V')(\A)\), \(r\in R(\A)\), set \(\Sigma\) of classes and \(\delta\in U(V')(F)\) of degree \(\ge0\) for \(d\); and if \(u\in R(\A)\), \(u-1\) has degree \(\ge2\) and \(\tilde w(u-1)\) has degree \(\ge1\), then
\[
  \mathcal Q^d_w(h,ur)=\frac{\psi_E\bigl(\frac12j_E\tr(\st^*\st\,u)\bigr)}{\psi_E\bigl(\frac12j_E\tr(\tilde w^*\tilde w\,u)\bigr)}\,\mathcal Q^d_w(h,r).
\]
The ratio is \(1\) for all such \(u\) exactly when \(\tilde w\) survives.

\begin{problem}\label{pr:covariance}
Let \(\gamma\in R(F)\) and \(\gamma'\in U(V')(F)\), and let \(\gamma'\) have degree \(\ge0\) for the weights \(d\) (no condition when \(d=0\)).  Then \(\mathcal Q^d_{\gamma'w\gamma^{-1}}(\gamma'h,\gamma r)=\mathcal Q^d_w(h,r)\).
\end{problem}

We use special cases of two facts on discrete groups.  Let \(G\) be a second-countable locally compact Hausdorff group with a left Haar measure, \(\Gamma\subset G\) a discrete subgroup acting on a countable set \(\mathcal X\), quotients by discrete subgroups carrying the induced measures, and \(\mathfrak k_{\mathfrak x}\) (\(\mathfrak x\in\mathcal X\)) measurable functions on \(G\) with \(\mathfrak k_{\gamma\mathfrak x}(\gamma g)=\mathfrak k_{\mathfrak x}(g)\).  If \(\sum_{\mathfrak x}\lvert\mathfrak k_{\mathfrak x}\rvert\) is integrable on \(\Gamma\backslash G\), then \(\mathfrak k_{\mathfrak x}\) is integrable on \(\Stab_\Gamma(\mathfrak x)\backslash G\) and \(\int_{\Gamma\backslash G}\sum_{\mathfrak x'\in\Gamma\mathfrak x}\mathfrak k_{\mathfrak x'}=\int_{\Stab_\Gamma(\mathfrak x)\backslash G}\mathfrak k_{\mathfrak x}\) \cite{Bou}.  And if \(\mathfrak f\) is continuous, integrable on \(\Gamma\backslash G\) and left \(\Gamma\)-invariant, with \(\mathfrak f(ug)=\zeta\mathfrak f(g)\) for all \(g\), one \(u\in G\) and one \(\zeta\ne1\) of modulus one, then \(\int_{\Gamma\backslash G}\mathfrak f=0\); Problem~\ref{pr:twisted-mean} is the finite case.  \emph{Status}: the first is classical; for the second we use the scalar, integrable case of the twisted-integral lemma in the proof.

\begin{problem}\label{pr:orbit-zero}
Let \(d\) be \(0\) or the grading of \(V'\), and let the class \(w\) fail the test for \(d\).  Then for every \(h\), \(\int_{[R]}\sum_{w'\in w R(F)}\mathcal Q^d_{w'}(h,r)\,dr=0\).
\end{problem}

Let \(w_+\) and \(\st_+\) be the matrices of the columns of positive grade of \(w\) and of \(\st\).  In the proof, if a rational \(w\) survives for \(d=0\), then \(w_+\) has linearly independent columns and \(\bar w_+^{t}J_{V'}w_+=\bar\st_+^{t}J_{V'}\st_+\); Problem~\ref{pr:hook-survivors} illustrates this.  Witt's theorem: if \(w_1,w_2\) are matrices over \(E\) of the same shape, with linearly independent columns and \(\bar w_1^{t}Jw_1=\bar w_2^{t}Jw_2\) for a nondegenerate hermitian \(J\), then \(gw_1=w_2\) for some invertible square \(g\) over \(E\) with \(\bar g^{t}Jg=J\) \cite{Sch}.  \emph{Status}: classical.

\begin{problem}\label{pr:witt}
Let \(w\) be a rational map with \(w_+\) of linearly independent columns and \(\bar w_+^{t}J_{V'}w_+=\bar\st_+^{t}J_{V'}\st_+\).  Then \(\gamma'w\) agrees with \(\st\) for some \(\gamma'\in U(V')(F)\).
\end{problem}

Let \(\Cc\) be the set of classes for the weights \(0\) that agree with \(\st\).  In the proof, if \(w\) agrees with \(\st\), \(\gamma\in R(\A)\) and \(\gamma'\in U(V')(\A)\), then \(\gamma'w\gamma\) agrees with \(\st\) if and only if \(\gamma'\in R'(\A)\); in particular \(\Cc\) is stable under \(w\mapsto w\gamma\), \(\gamma\in R(F)\).  Problem~\ref{pr:hook-orbit} illustrates this.

\begin{problem}\label{pr:block}
Let a group \(H\) act on a set \(\mathcal X\), let \(S\subseteq H\) be a subgroup and \(\Cc\subseteq\mathcal X\) a subset such that for \(\mathfrak x\in\Cc\), \(g\mathfrak x\in\Cc\) if and only if \(g\in S\).  Then \((q,\mathfrak x)\mapsto q_\rep\mathfrak x\), for any choice of representatives \(q_\rep\), is a bijection from \((H/S)\times\Cc\) onto the union of the translates of \(\Cc\).
\end{problem}

The second selection takes place inside the period: in the proof, for every \(h\),
\[
  \int_{[R]}\sum_{w\in\Cc}\mathcal Q^0_w(h,r)\,dr=\Wc'(h).
\]
For each datum, the sum of the constant terms over \(\Cc\) equals the sum over the classes agreeing with \(\st\) for the grading of \(V'\); every surviving representative of one of those lies in the class of \(\st\); \(\st\) survives, and \(R(F)\) fixes its class.  Problem~\ref{pr:hook-second} illustrates this.

\begin{problem}\label{pr:period-sum}
For every \(h\in U(V')(\A)\), the series on the right converges absolutely and
\[
  \int_{[R]}\Wc(r)\,\theta\bigl(\omega(r,h)\phi\bigr)\,dr=\sum_{q\in U(V')(F)/R'(F)}\Wc'(q_\rep^{-1}h).
\]
\end{problem}

\begin{problem}\label{pr:unfold}
Let \(\Gamma'\subseteq\Gamma\) be discrete subgroups of \(G\) as above, \(\mathfrak f\) a continuous left \(\Gamma'\)-invariant function and \(z\) a continuous left \(\Gamma\)-invariant function on \(G\), with \(\sum_{q\in\Gamma/\Gamma'}\lvert\mathfrak f(q_\rep^{-1}h)z(h)\rvert\) integrable on \(\Gamma\backslash G\).  Then \(\mathfrak fz\) is integrable on \(\Gamma'\backslash G\) and
\[
  \int_{\Gamma\backslash G}\sum_{q\in\Gamma/\Gamma'}\bigl(\mathfrak f(q_\rep^{-1}h)\,z(h)\bigr)\,dh=\int_{\Gamma'\backslash G}\mathfrak f(h)z(h)\,dh .
\]
\end{problem}

The height of \(h\in U(V')(\A)\) is the reciprocal of the first minimum (the length of a shortest nonzero lattice vector), cut off at \(1\), of the lattice of the space moved by \(h\).  In the proof, for every \(g_0\) and every large integer \(e\), each \(h\) has a \(\gamma'\in U(V')(F)\) with \(\lVert\omega(g_0,\gamma'h)\rVert\le2^e\) times the \(e\)-th power of the height of \(h\); we use \(g_0=1\).  This growth estimate is proved by reduction theory; compare \cite{KR}.  A function \(z\) on \(U(V')(\A)\) is \emph{automorphic} if it is continuous, left invariant under \(U(V')(F)\) and invariant under right translation by a neighbourhood of \(1\) in the finite adelic group, and \emph{rapidly decreasing} if for every integer \(b\ge0\), \(\lvert z(h)\rvert\) times the \(b\)-th power of the height of \(h\) is bounded.  For each rapidly decreasing \(z\), in the proof, \(\sum_{w\in\Sigma}\lvert\mathcal Q^0_w(h,r)z(h)\rvert\le C'\) for one constant \(C'\), every \(h\), \(r\) and every set \(\Sigma\) of classes, and \(\sum_q\lvert\Wc'(q_\rep^{-1}h)z(h)\rvert\le C'\operatorname{vol}[R]\), and \([U(V')]\) has finite volume.  Problems~\ref{pr:covariance}--\ref{pr:unfold}, with the facts stated in this section, prove the following.

\begin{theorem}\label{thm:transfer}
With \(E\), the frames, \(\omega\) and \(\Sc\) as in this section, let \(V\), \(V'\) be hermitian spaces over \(E\), each the orthogonal sum of a definite part and a string for the same \(N\), with \(\len V'=\len V-1\) or \(\len V=\len V'=0\); let \(R\), \(R'\) be their Bessel subgroups and \(\st\) the string map.  Let \(\phi\in\Sc\), \(\Wc\) be as in Definition~\ref{def:transfer}, and \(z\) be an automorphic, rapidly decreasing function on \(U(V')(\A)\).  Then \(\Wc'z\) is integrable on \(R'(F)\backslash U(V')(\A)\) and
\[
  \int_{[R]}\Wc(r)\int_{[U(V')]}\theta\bigl(\omega(r,h)\phi\bigr)\,z(h)\,dh\,dr=\int_{R'(F)\backslash U(V')(\A)}\Wc'(h)\,z(h)\,dh .
\]
\end{theorem}

\paragraph{What has been proved.}  A Bessel period of a theta function, paired with a form on \(V'\), is the integral of the transfer against the form over \(R'(F)\backslash U(V')(\A)\).  Used as given: the facts on \(\omega\), \(\Heis\) and \(\theta\) stated at the start of Section~\ref{sec:transfer}, the expansion of \(\theta\), the statements marked ``in the proof'', Witt's theorem and the facts on discrete groups.  If \(c_V=c_{V'}\), then in the proof, for each \(h\), the function \(r'\mapsto\Wc'(r'h)\) on \(R'(\A)\) satisfies the hypothesis on \(\Wc\) for \(V'\): the transfer is again a Bessel function, on a string of length \(\max(\len V-1,0)\).  Along the lower chain \(Y_0,\dots,Y_m\) the string of \(Y_j\) has length \(m-1-j\) for \(j<m\), and that of \(Y_m\) has length \(0\); in the proof the gradings of the \(Y_j\), \(j<m\), have one common constant, and that of \(Y_m\) has constant \(1\).

\begin{leanremark}
{\raggedright Theorem~\ref{thm:transfer} is \tl{New.Grading.integral\_besselGroup\_unfold} in the Lean proof \cite{GL}.  The adelic definitions specialize \tl{New.Grading}, \tl{stringMap}, \tl{constTerm}, \tl{unfoldStep}, \tl{Survives} and \tl{HasDegreeGe}; the finite examples use the same matrix formulas.  The expansion of \(\theta\) is \tl{hasSum\_constTerm}, the matrix form of Witt's theorem is \tl{New.witt}, and the growth estimate is \tl{eventually\_exists\_norm\_omega\_le\_mul\_height\_pow}.  The statements marked ``in the proof'' are at these places of \cite{GL}, commit \texttt{72fead1}: \glsrc{Lib/Weil/Periods/Grading.lean}{134-L184}{the descriptions of \(\Uc\) and \(R\)}, \glsrc{Lib/Weil/Periods/BesselIntegral.lean}{133}{compactness of \([R]\)}, \glsrc{Lib/Weil/Periods/Equivariant.lean}{215-L238}{the character \(\psi_\Uc\)} and \glsrc{Lib/Weil/Periods/Grading.lean}{875}{\(\st^*\st\)}, \glsrc{Lib/Weil/Periods/Grading.lean}{809}{the degree of \(\st\)}, \glsrc{Lib/Weil/Periods/ConstantTerm.lean}{419-L452}{class independence of \(\CT^d_w\)}, \glsrc{Lib/Weil/Periods/ConstantTerm.lean}{852-L880}{continuity and invariance of \(\Wc'\)}, \glsrc{Lib/Weil/Periods/Unfolding.lean}{136-L242}{the terms \(\mathcal Q^d_w\)} and \glsrc{Lib/Weil/Periods/ConstantTerm.lean}{800-L823}{their phase law}, \glsrc{Lib/Weil/Periods/BesselIntegral.lean}{722-L821}{the Gram condition}, \glsrc{Lib/Weil/Periods/BesselIntegral.lean}{645-L678}{preservation of agreement}, \glsrc{Lib/Weil/Periods/Unfolding.lean}{370-L398}{the second selection}, \glsrc{Lib/Weil/Periods/Unfolding.lean}{459-L510}{the two estimates} and \glsrc{Lib/Adelic/Reduction/FiniteVolume.lean}{772}{the finite volume}, \glsrc{Lib/Weil/Periods/ConstantTerm.lean}{887-L912}{the case \(c_V=c_{V'}\)}, \glsrc{New/UpperBound/Seesaw/EdgePairEvaluation.lean}{1140-L1162}{the chain gradings}, \glsrc{New/UpperBound/Seesaw/SeesawIdentity.lean}{197-L223}{the transposition}, \glsrc{Lib/Adelic/QuotientMeasure/IwasawaUnfolding.lean}{302-L316}{the orbit fact} and \glsrc{Lib/Adelic/QuotientMeasure/UnfoldingInStages.lean}{362-L393}{the twisted integral}.  See also the \glsrc{New/UpperBound/Seesaw/UnfoldingStages.lean}{688-L705}{period at a chain step}, the \glsrc{Lib/Weil/Periods/Unfolding.lean}{346-L366}{selection at zero weights}, and the \glsrc{Lib/Weil/Periods/Unfolding.lean}{435-L454}{sum over the translates}.  The conventions are those for \glsrc{Lib/Adelic/Unitary/Defs.lean}{498-L506}{the adelic frame}, \glsrc{Lib/Adelic/Adeles/AdditiveCharacter.lean}{153-L160}{the standard additive character}, \glsrc{Lib/Weil/Construction/WeilData.lean}{69-L89}{the Weil data}, \glsrc{Lib/ClassField/HeckeCharacter/Basic.lean}{61-L74}{the norm-one characters}, and \glsrc{Lib/Weil/Construction/AdelicRepresentation.lean}{241-L250}{the space of data}.  For the action and the trace phase see \glsrc{Lib/Weil/DualPair/WeilAction.lean}{761-L777}{the twisted action and its bicharacter}, \glsrc{Lib/Weil/Periods/ConstantTerm.lean}{102-L111}{Heisenberg covariance}, \glsrc{Lib/Weil/DualPair/WeilAction.lean}{616-L623}{rational invariance of the theta kernel}, and \glsrc{Lib/Weil/Periods/ConstantTerm.lean}{423-L435}{rational Heisenberg translation}.  The function hypotheses are \glsrc{Lib/Weil/Periods/Equivariant.lean}{273-L284}{Bessel equivariance}, \glsrc{Lib/Adelic/AutomorphicForms/Basic.lean}{102-L111}{automorphy}, and \glsrc{Lib/Adelic/Reduction/FirstMinimumShortBases.lean}{281-L298}{the height and rapid decay}.\par}
\end{leanremark}


\section*{Hints}

\subsection*{\S\ref{sec:plane}}
\begin{description}
\item[\ref{pr:coefficient}] Exchange the two sums and use \(\sum_b\zeta_3^{b(Q(v)-1)}\).
\item[\ref{pr:plane-orbit}] \(-1\) is not a square in \(\F_3\), so \(Q\) is anisotropic.
\item[\ref{pr:plane-unfold}] Problems~\ref{pr:coefficient} and~\ref{pr:plane-orbit}; substitute \(h\mapsto\gamma h\), as in Problem~\ref{bp-pr:unfold}.
\end{description}

\subsection*{\S\ref{sec:strings}}
\begin{description}
\item[\ref{pr:hook-unipotent}] Expand \(\bar u^{t}J_Vu=J_V\) entry by entry.
\item[\ref{pr:hook-phase}] Compute the shift \(\st^*\st\) on the string basis.
\end{description}

\subsection*{\S\ref{sec:survival}}
\begin{description}
\item[\ref{pr:hook-survivors}] Which \(u\) are admissible for \(v=0\), and which for \(v\ne0\)?  Then Problem~\ref{pr:hook-phase}.
\item[\ref{pr:hook-orbit}] Count first, then second columns.  The Bessel subgroup of a string of length two fixes \(f'_0\).
\item[\ref{pr:hook-second}] Compare the entries used for agreement with those of negative degree; then compute the two phases.
\end{description}

\subsection*{\S\ref{sec:transfer}}
\begin{description}
\item[\ref{pr:covariance}] Move \(\omega(\gamma,\gamma')\) to the left through the Heisenberg operators, then substitute \(y\mapsto\gamma'y\gamma^{-1}\) in the mean.  An element of \(R(\A)\) preserves the span of the basis vectors of grade \(\ge i\) for each \(i\), so it does not lower grades (use unitarity for the grades \(<0\) and for the definite part).
\item[\ref{pr:orbit-zero}] Unfold the orbit with \(\Gamma=R(F)\) and Problem~\ref{pr:covariance}; on the stabilizer quotient use the vanishing of twisted integrals.
\item[\ref{pr:witt}] Witt's theorem in \(V'\), for the columns of positive grade.
\item[\ref{pr:period-sum}] Expand \(\theta\) for the weights \(0\), and sort the classes into the translates of \(\Cc\) and the rest.
\item[\ref{pr:unfold}] The orbit fact for \(\Gamma\) acting on \(\Gamma/\Gamma'\); compare Problem~\ref{ps-pr:unfold}.
\end{description}

\begin{thebibliography}{Bou}
\bibitem[Bou]{Bou} N.~Bourbaki, \emph{Int\'egration}, Chapitres~7 et~8, Hermann, Paris, 1963.
\bibitem[GL]{GL} L.~Sernau, \emph{GRH for Dirichlet \(L\)-functions, in Lean~4}, a shorter rewriting of the Lean proof, text by Claude, \url{https://github.com/lsernau/grh-lean}, commit \texttt{72fead1}, readable edition \texttt{grh/}.
\bibitem[KR]{KR} S.~S.~Kudla and S.~Rallis, \emph{A regularized Siegel--Weil formula: the first term identity}, Ann.\ of Math.\ \textbf{140} (1994), 1--80.
\bibitem[Ku]{Ku} S.~S.~Kudla, \emph{Splitting metaplectic covers of dual reductive pairs}, Israel J.\ Math.\ \textbf{87} (1994), 361--401.
\bibitem[Sch]{Sch} W.~Scharlau, \emph{Quadratic and Hermitian Forms}, Grundlehren Math.\ Wiss.\ 270, Springer, 1985.
\bibitem[We]{We} A.~Weil, \emph{Sur certains groupes d'op\'erateurs unitaires}, Acta Math.\ \textbf{111} (1964), 143--211.
\end{thebibliography}

\end{document}
